% Ampliación aditiva: acción, área, memoria sectorial y termometría.
% Propietarios: 70_horizontes_e_informacion.tex; sección térmica antecedente;
% 71_energia_libre_holografica_cerrada.tex del corpus integral.
% Se mantienen las condiciones de la realización cosmológica declarada.

\subsection{Conversión térmica entre el calendario y el horizonte}
\label{viii:subsec:conversion-termica-calendario-horizonte}

La acción, el incremento de área y la memoria angular admiten una
composición energética explícita. Se conserva una misma orientación de
acción, la rama positiva de Barbero--Immirzi y la realización de
\eqref{vii:eq:horizonte-cosmologico}. Abreviamos
\[
 L=\ell_P=L_\alpha,\qquad
 E_P=\frac{h}{108t_0}=\frac{\hbar c}{L},\qquad
 q_{\mathcal A}=\delta\mathcal A=\gamma_{\rm BI}a_0L^2.
\]
El símbolo \(q_{\mathcal A}\) denota aquí un incremento de área.
El factor angular \(a_0\), el funcional \(\gamma_{\rm BI}\) y
la acción son las magnitudes obtenidas en las secciones precedentes.
Su dominio conserva la incidencia hexádico--octádica, la orientación y
el calendario del mismo estado enriquecido APP--TRIT--TPK.

\begin{proposition}[Composición de las escalas areal y térmica]
\label{viii:prop:composicion-area-termica}
Con la tensión de~\eqref{vii:eq:tension-informacion-area}, se cumple
\begin{equation}
 \frac{\mathcal T_{I/A}q_{\mathcal A}}{\log3}
 =\frac{E_P}{\log3}
 =\frac{h}{108t_0\log3}
 =k_B\Theta_{\rm clk}.
 \label{viii:eq:composicion-area-termica}
\end{equation}
Para \(R=\zeta L\), con \(\zeta>0\), sea
\(\tau_R=k_BT_{\rm cos}(R)\). Entonces
\begin{equation}
 \tau_R=\frac{E_P}{2\pi\zeta},\qquad
 \frac{T_{\rm cos}(R)}{\Theta_{\rm clk}}
 =\frac{\log3}{2\pi\zeta}.
 \label{viii:eq:conversion-dos-temperaturas}
\end{equation}
En particular, el cociente en \(R=L\) es \(\log3/(2\pi)\).
\end{proposition}
\begin{proof}
La definición \(\mathcal T_{I/A}=E_P/q_{\mathcal A}\) y la
normalización trítica de~\eqref{vii:eq:unidad-accion-informacion}
dan~\eqref{viii:eq:composicion-area-termica}. La expresión
gravitatoria de la misma tensión conduce al mismo resultado:
\[
 \frac{c^4}{\gamma_{\rm BI}a_0GL}
        \gamma_{\rm BI}a_0L^2
 =\frac{c^4L}{G}=\frac{\hbar c}{L}=E_P,
\]
donde se ha utilizado \(G\hbar=c^3L^2\).
Por~\eqref{vii:eq:horizonte-cosmologico},
\(\tau_R=\hbar c/(2\pi R)=E_P/(2\pi\zeta)\).
Dividir esta igualdad por \(k_B\Theta_{\rm clk}=E_P/\log3\)
prueba el cociente de temperaturas.
\end{proof}

En la misma sección de acción, la composición gravitatoria queda
expresada por
\begin{equation}
 Gk_B\Theta_{\rm clk}\log3=GE_P=c^4L,
 \qquad G\tau_R=\frac{c^4L}{2\pi\zeta}.
 \label{viii:eq:composicion-termica-gravitatoria}
\end{equation}
Ambas identidades siguen de \(G\hbar=c^3L^2\) y de las escalas
anteriores.

La tensión convierte el incremento areal en la energía de decisión.
La normalización trítica y la periodicidad cosmológica determinan dos
secciones térmicas de esa misma energía, relacionadas por
\eqref{viii:eq:conversion-dos-temperaturas}. La conversión conserva
ambas normalizaciones y la función de Barbero--Immirzi en la resolución
geométrica.

\subsection{Respuesta energética de la memoria sectorial}
\label{viii:subsec:respuesta-energetica-sectorial}

Fijemos \(R>0\) y el estado fiel
\(\rho_0=\rho_\partial(\Delta_0)\), con \(\Delta_0>0\),
de~\eqref{vii:eq:estado-sectorial}. Su operador modular es
\(\mathsf K_0=(\log Z_0)I+\Delta_0\mathsf D_\partial\).
La realización de horizonte asigna a cada estado de densidad
\(\sigma\) las magnitudes de respuesta
\begin{equation}
\begin{aligned}
 \mathcal E_R(\sigma)
 &=E_\Lambda(R)+\tau_R
   \operatorname{Tr}[(\sigma-\rho_0)\mathsf K_0],\\
 \mathcal S_R(\sigma)
 &=S_{\rm H}(R)+k_B[s(\sigma)-s(\rho_0)].
\end{aligned}
 \label{viii:eq:respuesta-horizonte-definida}
\end{equation}
Se trata de la respuesta del estado modular en la realización geométrica
ya especificada. El radio y la acción fijan su escala energética
\(\tau_R\); las ocupaciones conservan la procedencia de cada canal.

\Needspace{28\baselineskip}
\begin{theorem}[Hamiltoniano de respuesta y balance informacional]
\label{viii:thm:hamiltoniano-respuesta-termica}
El operador de diferencias energéticas correspondiente a
\eqref{viii:eq:respuesta-horizonte-definida} puede elegirse como
\begin{equation}
 \mathsf H_R=\tau_R\Delta_0\mathsf D_\partial.
 \label{viii:eq:hamiltoniano-respuesta}
\end{equation}
Su estado de Gibbs a escala energética \(\tau_R\) es \(\rho_0\).
Para todo estado de densidad \(\sigma\) en el espacio sectorial,
\begin{equation}
 \mathcal E_R(\sigma)-T_{\rm cos}(R)\mathcal S_R(\sigma)
 =\tau_R\mathcal D(\sigma\Vert\rho_0)\geq0.
 \label{viii:eq:balance-libre-termico}
\end{equation}
En una variación diferenciable de estados alrededor de \(\rho_0\),
con \(R\) y \(\mathsf H_R\) fijos, se obtiene
\begin{equation}
 \delta\mathcal E_R
 =\tau_R\,\delta s
 =T_{\rm cos}(R)\,\delta\mathcal S_R.
 \label{viii:eq:primera-ley-energetica}
\end{equation}
\end{theorem}
\begin{proof}
Las trazas de \(\sigma\) y \(\rho_0\) son uno, por lo que la
diferencia elimina el término escalar \(\log Z_0\). De aquí resulta
\(\mathcal E_R(\sigma)-E_\Lambda(R)
=\operatorname{Tr}[(\sigma-\rho_0)\mathsf H_R]\).
Además,
\[
 \frac{e^{-\mathsf H_R/\tau_R}}
 {\operatorname{Tr}e^{-\mathsf H_R/\tau_R}}
 =\frac{e^{-\Delta_0\mathsf D_\partial}}{Z_0}=\rho_0.
\]
La igualdad \(E_\Lambda=T_{\rm cos}S_{\rm H}\) cancela la
referencia geométrica en el balance. La identidad finita
\eqref{vii:eq:ley-modular-finita}, escrita como
\(s(\sigma)-s(\rho_0)
=\operatorname{Tr}[(\sigma-\rho_0)\mathsf K_0]
-\mathcal D(\sigma\Vert\rho_0)\), demuestra
\eqref{viii:eq:balance-libre-termico}. Su término tangente es la
primera ley modular ya probada y produce
\eqref{viii:eq:primera-ley-energetica}.
\end{proof}

Para un incremento unitario de ocupación, los saltos energéticos son
\(\epsilon_{90}=90\tau_R\Delta_0\) y
\(\epsilon_{120}=120\tau_R\Delta_0\). En ambos sectores,
\begin{equation}
 \frac{\epsilon_m}{m\Delta_0}=\tau_R,
 \qquad m\in\{90,120\}.
 \label{viii:eq:pendiente-termica-comun}
\end{equation}
Así, el cociente \(4/3\) entre pesos areales registra los diferentes
saltos del mismo Hamiltoniano. El cociente entre un salto energético y
la variación correspondiente de \(-\log p\) conserva una escala común.
La normalización del estado elimina sólo una constante escalar.

\begin{corollary}[Susceptibilidad energética a Hamiltoniano fijo]
\label{viii:cor:respuesta-escala-termica}
Para \(\rho_\tau=e^{-\mathsf H_R/\tau}/
\operatorname{Tr}e^{-\mathsf H_R/\tau}\), con \(\tau>0\),
\begin{equation}
 \frac{\mathrm d\langle\mathsf H_R\rangle_\tau}{\mathrm d\tau}
 =\frac{\operatorname{Var}_{\rho_\tau}(\mathsf H_R)}{\tau^2}>0.
 \label{viii:eq:susceptibilidad-termica}
\end{equation}
En \(\tau=\tau_R\), el resultado es
\(\Delta_0^2\operatorname{Var}_{\rho_0}(\mathsf D_\partial)\).
\end{corollary}
\begin{proof}
La derivación de la suma finita de exponenciales da la varianza dividida
por \(\tau^2\). El estado es fiel y \(\mathsf H_R\) tiene más de
un autovalor, de modo que la varianza es positiva. La sustitución de
\eqref{viii:eq:hamiltoniano-respuesta} da la evaluación indicada.
\end{proof}

Esta derivada varía la escala de Gibbs y mantiene fijo el Hamiltoniano.
La variación radial de~\eqref{vii:eq:diferencial-horizonte} modifica
también \(\mathsf H_R\) y conserva sus términos geométricos propios.

\subsection{Sección termométrica y conservación del registro}
\label{viii:subsec:seccion-termometrica-registro}

El estado \(\rho_0\) y el Hamiltoniano de respuesta permiten fijar una
coordenada termométrica anterior a cualquier lector numérico candidato.
Para el radio fijado, sea \(\mathcal L_\Theta=\mathbb R u_R\) una
recta orientada cuya sección positiva \(u_R\) corresponde al estado
de referencia. En los estados fieles de Gibbs del mismo Hamiltoniano,
\begin{equation}
 \rho_\beta=\frac{e^{-\beta\mathsf H_R}}
                  {\operatorname{Tr}e^{-\beta\mathsf H_R}},\qquad
 \Theta_R(\rho_\beta)=\frac{1}{\beta\tau_R}u_R,
 \quad \beta>0.
 \label{viii:eq:lector-termometrico}
\end{equation}
La identidad operatoria
\(-\log\rho_\beta=\beta\mathsf H_R+(\log Z_\beta)I\)
determina \(\beta\): restar dos autovalores con distinta energía
da su cociente único. En particular,
\(\Theta_R(\rho_0)=u_R\).
La recta \(\mathcal L_E\) es el espacio unidimensional de las
magnitudes energéticas.

\Needspace{23\baselineskip}
\begin{theorem}[Transductor positivo y criterio de identificación]
\label{viii:thm:identificacion-transductor}
Existe un único mapa lineal positivo
\(B_R:\mathcal L_\Theta\to\mathcal L_E\) con
\(B_R(u_R)=\tau_R\). Satisface
\begin{equation}
 B_R(\Theta_R(\rho_\beta))=\beta^{-1}.
 \label{viii:eq:transductor-gibbs}
\end{equation}
La sección cronológica compatible es
\begin{equation}
 \Theta_{{\rm clk},R}
 =\frac{E_P}{\tau_R\log3}u_R
 =\frac{2\pi R}{L\log3}u_R,
 \qquad
 B_R(\Theta_{{\rm clk},R})=
 \frac{\mathcal T_{I/A}q_{\mathcal A}}{\log3}.
 \label{viii:eq:seccion-cronologica-compatible}
\end{equation}
Para cualquier otro mapa lineal \(B_*\) entre las mismas rectas,
\begin{equation}
 B_*=B_R
 \quad\Longleftrightarrow\quad
 B_*(\Theta_{{\rm clk},R})=\frac{h}{108t_0\log3}.
 \label{viii:eq:criterio-una-seccion}
\end{equation}
\end{theorem}
\begin{proof}
La sección \(u_R\) es una base de una recta, de modo que prescribir
su imagen define un único mapa lineal; la positividad de \(\tau_R\)
asegura la positividad del mapa. Sustituir
\eqref{viii:eq:lector-termometrico} prueba
\eqref{viii:eq:transductor-gibbs}. La expresión de \(\tau_R\) y
\eqref{viii:eq:composicion-area-termica} dan
\eqref{viii:eq:seccion-cronologica-compatible}. Finalmente,
\(\Theta_{{\rm clk},R}\) es no nula. Dos mapas lineales de una
recta coinciden si y sólo si coinciden sobre esa sección.
\end{proof}

El criterio compara lectores en una sección obtenida previamente del
estado, la respuesta energética y la normalización trítica. La aplicación
a una coordenada numérica concreta requiere evaluar ese lector en la
sección especificada. Las bases dimensionales se transportan conjuntamente
con los mapas y con las coordenadas térmicas que expresan.

\begin{proposition}[Transporte de la respuesta con memoria]
\label{viii:prop:transporte-termico-memoria}
Sea \(J\) una isometría que conserva las etiquetas completas del
registro, y considérese su imagen como espacio receptor. Para
\(\mathsf H'_R=J\mathsf H_RJ^*\), \(\rho'_0=J\rho_0J^*\)
y \(\sigma'=J\sigma J^*\), se conservan las expectativas energéticas,
las entropías y el balance~\eqref{viii:eq:balance-libre-termico}.
También se cumple \(J\mathsf H_R=\mathsf H'_RJ\).
\end{proposition}
\begin{proof}
Sobre la imagen, \(J\) es unitario y transporta el cálculo funcional.
La ciclicidad de la traza conserva expectativas; los mismos autovalores
no nulos conservan entropía y entropía relativa. La identidad de
entrelazamiento usa \(J^*J=I\). Sustituir estas igualdades en el balance
prueba su conservación.
\end{proof}

La misma equivalencia admite otra escala \(\tau'=a\tau_R\),
con \(a>0\), y otro origen energético \(e'\). Si
\(\mathsf H'=e'I+aJ\mathsf H_RJ^*\), entonces
\[
 \frac{\mathsf H'-e'I}{\tau'}
 =J\frac{\mathsf H_R}{\tau_R}J^*,\qquad
 \frac{e^{-\mathsf H'/\tau'}}
      {\operatorname{Tr}_{\operatorname{im}J}e^{-\mathsf H'/\tau'}}
 =J\rho_0J^*.
\]
El factor escalar \(e^{-e'/\tau'}\) se cancela en el estado;
las diferencias energéticas se multiplican por \(a\), y la entropía
relativa permanece invariante. El balance correspondiente tiene, por
tanto, coeficiente \(\tau'\). La expectativa energética completa
conserva el origen: \(\langle\mathsf H'\rangle
=e'+a\langle\mathsf H_R\rangle\). La comparación de realizaciones
se expresa así mediante sus exponentes centrados y sus escalas propias.
Si la expectativa alimenta las lecturas \(m=E/c^2\) y
\(r_g=GE/c^4\), mantener \(c,G\) da
\[
 m'=\frac{e'}{c^2}+am,\qquad
 r'_g=\frac{Ge'}{c^4}+ar_g.
\]
La normalización estadística y el transporte de magnitudes mantienen,
por tanto, funciones distintas para el mismo origen registrado.

La composición enlaza el incremento areal de Barbero--Immirzi, la energía
del calendario y la distribución angular del borde. Su transporte
retiene la procedencia sectorial y el archivo: el retorno de fase
continúa acompañado por el avance de memoria del estado enriquecido.
